% optsets-proofs.tex -- nonunique optimal sets and exact boundary output.
% Fragment without preamble.  Uses main-text labels: lem:round, prop:filter,
% eq:mesh, eq:messages, eq:dpwork, lem:contraction, lem:localize,
% eq:witnessradius, eq:stagegap, lem:states, eq:statecap, eq:logabsorb,
% thm:approx, sec:rate, sec:model, sec:polynomial, thm:generalfinite.
% New bibliography keys (see report): QiuYildirim2023, JeyakumarRubinovWu2006.

\subsection{Nonunique optimal sets}\label{sec:optsets}

The point-growth analysis of Section~\ref{sec:rate} localizes every retained
coordinate near one optimizer. When the optimal set is not a single point,
its coordinate projections can fill whole intervals, and that localization
fails. This section answers three questions. Which compact descriptions of
the whole optimal set can be computed with the decomposition? Why can no
algorithm that uses only function values give accuracy-independent work under
growth toward a set? And for which class can an unknown optimal set be found
exactly, with the growth constant unknown?

\paragraph{Conventions.}
In this section
\[
F(x)=\tfrac12x^{\mathsf T}Hx+b^{\mathsf T}x+c
\]
with rational symmetric $H$, rational $b,c$, and $X=\prod_iX_i$ as in
Section~\ref{sec:model}, preprocessed so that $l_i<u_i$ for every $i$. Write
$\hat X=\prod_i[l_i,u_i]$ for the continuous hull, $S=\arg\min_XF$ and
$\OPT=\min_XF$. The interaction graph has an edge $\{i,j\}$ when $H_{ij}\ne0$.
The supplied tree decomposition $(T,(B_t))$ covers it, and each monomial is
assigned to one bag that contains its variables. A number $g>0$ is a
\emph{set-growth constant} if
\begin{equation}\label{eq:setgrowth}
F(x)-\OPT\ge g\dist(x,S)^2\qquad(x\in X),
\end{equation}
and then $\kappa=\max\{1,L/g\}$. For a quadratic, every rational
$L\ge\max_i\max\{H_{ii},0\}$ with $L>0$ is a valid coordinate-curvature bound.
For a bounded continuous box, some set-growth constant always exists: apply
the error bound of Luo and Sturm \cite[Theorem~3.3]{LuoSturm2000} to the
quadratic $F-\OPT$ on the polytope $X$. Its size is not controlled, and
nothing below computes or certifies it. $\mathrm{Id}$ denotes an identity
matrix.

\subsubsection{Coordinatewise concave quadratics: every optimizer at once}

Write $E_i=\{l_i,u_i\}$ and $E(B)=\prod_{i\in B}E_i$. For $x\in\hat X$ and
$v_i\in E_i$ put
\[
e_i(v_i;x_i)=\begin{cases}u_i-x_i,&v_i=l_i,\\ x_i-l_i,&v_i=u_i,\end{cases}
\qquad
\pi_t(v;x)=\prod_{i\in B_t}\frac{e_i(v_i;x_i)}{u_i-l_i}\quad(v\in E(B_t)).
\]
Thus $\pi_t(\cdot;x)$ is the law of $Y_{B_t}$ when every coordinate of $x$
is rounded independently to an endpoint with mean $x_i$. Root $T$ at $t_0$,
let $\mathrm{ch}(t)$ be the children of $t$, and write
$B_t^{\uparrow}=B_t\cap B_{\mathrm{pa}(t)}$, with $B_{t_0}^{\uparrow}=\emptyset$.
Let $\phi_t$ be the sum of the monomials assigned to $t$.

\begin{lemma}[Endpoint Bellman residual identity]\label{lem:endpointid}
Assume $H_{ii}\le0$ for every $i$. Compute, in post-order and with all
arguments ranging over endpoint labels,
\begin{align*}
M_t(v_{B_t^\uparrow})&=\min_{v_{B_t\setminus B_t^\uparrow}}
\Bigl[\phi_t(v_{B_t})+\sum_{t'\in\mathrm{ch}(t)}M_{t'}(v_{B_{t'}^\uparrow})\Bigr],\\
r_t(v_{B_t})&=\phi_t(v_{B_t})+\sum_{t'\in\mathrm{ch}(t)}M_{t'}(v_{B_{t'}^\uparrow})
-M_t(v_{B_t^\uparrow})\ \ge0 .
\end{align*}
Then the scalar $M_{t_0}$ equals $\OPT$, it is attained at an endpoint
assignment, and for every $x\in\hat X$
\begin{equation}\label{eq:endpointid}
F(x)-\OPT=\sum_{t}\sum_{v\in E(B_t)}r_t(v)\,\pi_t(v;x)
+\frac12\sum_{i=1}^n(-H_{ii})(x_i-l_i)(u_i-x_i).
\end{equation}
All tables and residuals are computed in
$O\bigl(p\,2^p(N+|\mathcal A|)\bigr)$ table operations on rationals of
bit length polynomial in $I$.
\end{lemma}

\begin{proof}
Since $B_{t'}^{\uparrow}\subseteq B_t$ for a child $t'$, the bracket is a
function of $v_{B_t}$, and $r_t\ge0$ because $M_t$ is its minimum over the
labels of $B_t\setminus B_t^\uparrow$.

\emph{Telescoping.} Let $v\in E([n])$. Every nonroot table $M_t$ occurs in
$\sum_tr_t(v_{B_t})$ once with a plus sign, in the residual of its parent, and
once with a minus sign. Hence $\sum_tr_t(v_{B_t})=\sum_t\phi_t(v_{B_t})-M_{t_0}
=F(v)-M_{t_0}$.

\emph{Expectation.} Fix $x\in\hat X$ and let $Y$ be the independent endpoint
rounding with $\Pr(Y_i=v_i)=e_i(v_i;x_i)/(u_i-l_i)$. Then $\E Y_i=x_i$,
$\operatorname{Var}Y_i=(x_i-l_i)(u_i-x_i)$, $\E Y_iY_j=x_ix_j$ for $i\ne j$,
and $\E Y_i^2=x_i^2+\operatorname{Var}Y_i$. Therefore
$\E F(Y)=F(x)+\frac12\sum_iH_{ii}(x_i-l_i)(u_i-x_i)$, while
$\E r_t(Y_{B_t})=\sum_vr_t(v)\pi_t(v;x)$. Taking expectations in the
telescoping identity gives \eqref{eq:endpointid} with $M_{t_0}$ in place
of $\OPT$.

\emph{Value.} Every term on the right of \eqref{eq:endpointid} is
nonnegative on $\hat X$, so $F\ge M_{t_0}$ on $\hat X\supseteq X$. Trace back:
choose labels of $B_{t_0}$ attaining $M_{t_0}$; for each child $t'$ of a
processed bag $t$, choose the labels of $B_{t'}\setminus B_{t'}^\uparrow$
attaining $M_{t'}(v_{B_{t'}^\uparrow})$. By running intersection, a variable
belongs to $B_t\setminus B_t^\uparrow$ exactly for the bag $t$ nearest the
root among those containing it, and to $B_{t}^{\uparrow}$ for every other bag
containing it. Hence every variable receives exactly one label, and the
resulting endpoint assignment $v^*$ has all residuals zero. By telescoping
$F(v^*)=M_{t_0}$. Both endpoints of every $X_i$ are feasible, so $v^*\in X$
and $\OPT=M_{t_0}$.

\emph{Checking.} The proof uses only the nonnegativity of the residuals and
one endpoint assignment with zero residuals. A checker therefore needs the
tables $M_t$ and that assignment, not the minimization that produced them; any
tables with these two properties give \eqref{eq:endpointid}.

\emph{Cost.} A bag table has at most $2^p$ entries, and children are combined
by addition. By induction on the tree, every message entry is the sum of the
assigned monomials over the subtree at one endpoint assignment. Its bit length
is therefore bounded by the bit length of a common denominator of all endpoint
monomial values plus $\log_2$ of the number of monomials times their maximum
magnitude, both polynomial in $I$.
\end{proof}

\begin{theorem}[All optimizers of a coordinatewise concave quadratic]\label{thm:endpointset}
Under the hypothesis of Lemma~\ref{lem:endpointid}, a point $x\in X$ is a
global optimizer if and only if
\begin{enumerate}
\item[(a)] $x_i\in\{l_i,u_i\}$ for every $i$ with $H_{ii}<0$, and
\item[(b)] $\prod_{i\in B_t}e_i(v_i;x_i)=0$ for every bag $t$ and every
$v\in E(B_t)$ with $r_t(v)>0$.
\end{enumerate}
These are at most $n+N2^p$ factored polynomial equations of degree at most
$p$, together with the original bounds and integrality. They are computed and
checked in $2^{O(p)}\poly(I)$ bit operations, without a growth assumption and
without enumerating interior integer labels or optimal components.
\end{theorem}

\begin{proof}
By \eqref{eq:endpointid}, $x$ is optimal exactly when every nonnegative term
vanishes. The diagonal terms vanish if and only if (a) holds, and
$\pi_t(v;x)=0$ exactly when one factor $e_i(v_i;x_i)$ is zero.
\end{proof}

Condition (b) is a support condition, not a condition on coordinate
projections. For $F=x+y-2xy=x(1-y)+y(1-x)$ on $[0,1]^2$, both optimal
corners $(0,0)$ and $(1,1)$ use every endpoint of every coordinate, yet the
two positive residuals give $x(1-y)=0$ and $y(1-x)=0$, which exclude every
other point. A sum of $m$ disjoint copies has $2^m$ isolated optimizers and
a description with $2m$ equations.

The equations have a finite combinatorial form. For $x\in\hat X$ define its
\emph{face pattern} $\sigma(x)\in\{\mathsf L,\mathsf U,\ast\}^n$ by
$\sigma_i=\mathsf L$ if $x_i=l_i$, $\sigma_i=\mathsf U$ if $x_i=u_i$, and
$\sigma_i=\ast$ otherwise. For a pattern $\tau$ on $B\subseteq[n]$ let
$\operatorname{cube}(\tau)\subseteq E(B)$ consist of the $v$ with $v_i=l_i$
where $\tau_i=\mathsf L$ and $v_i=u_i$ where $\tau_i=\mathsf U$. Then
$\pi_t(v;x)>0$ if and only if $v\in\operatorname{cube}(\sigma(x)_{B_t})$.
Let $\mathcal P_i=\{\mathsf L,\mathsf U\}$, with $\ast$ added when $H_{ii}=0$
and $X_i$ has a point strictly between $l_i$ and $u_i$. Call
$\sigma\in\prod_i\mathcal P_i$ \emph{admissible} if, for every bag,
$\operatorname{cube}(\sigma_{B_t})\subseteq Z_t:=\{v\in E(B_t):r_t(v)=0\}$.

\begin{corollary}[Optimal faces form a tree-structured constraint problem]\label{cor:facecsp}
Under the hypothesis of Lemma~\ref{lem:endpointid}:
\begin{enumerate}
\item[(i)] $x\in X$ is optimal if and only if $\sigma(x)\in\prod_i\mathcal P_i$
and $\sigma(x)$ is admissible.
\item[(ii)] Replacing any $\ast$ of an admissible pattern by $\mathsf L$ or
$\mathsf U$ keeps it admissible, and
$S=\bigcup_{\sigma\text{ admissible}}\bigl(X\cap\bar F(\sigma)\bigr)$, where
$\bar F(\sigma)=\{x:x_i=l_i\text{ if }\sigma_i=\mathsf L,\ x_i=u_i\text{ if }
\sigma_i=\mathsf U\}$ is a closed face of $\hat X$. Distinct admissible
patterns give disjoint nonempty sets $\{x\in X:\sigma(x)=\sigma\}$.
\item[(iii)] Admissible patterns are the solutions of a constraint problem on
the same tree decomposition, with at most three labels per variable and bag
relations
$\mathcal A_t=\{\tau\in\prod_{i\in B_t}\mathcal P_i:\operatorname{cube}(\tau)
\subseteq Z_t\}$. Consequently, in $2^{O(p)}(N+n)\poly(I)$ bit operations
one can decide whether the optimizer is unique, compute the dimension of $S$,
count the admissible patterns, count $|S|$ when $S$ is finite, compute
$\dist(y,S)^2$ and a nearest optimizer for a rational $y$, and minimize a
linear function over $S$.
\end{enumerate}
\end{corollary}

\begin{proof}
(i) Condition (a) of Theorem~\ref{thm:endpointset} says $\sigma_i(x)\ne\ast$
when $H_{ii}<0$; and $\sigma_i(x)=\ast$ is possible only if $X_i$ has an
interior point. Condition (b) says that every $v$ with $\pi_t(v;x)>0$ has
$r_t(v)=0$, that is, $\operatorname{cube}(\sigma(x)_{B_t})\subseteq Z_t$.

(ii) Replacing $\ast$ by an endpoint shrinks every cube. If $\sigma$ is
admissible and $x\in X\cap\bar F(\sigma)$, then $\sigma(x)$ is obtained from
$\sigma$ by replacing some $\ast$ entries by endpoints, so $x$ is optimal by
(i). Conversely $x\in S$ lies in $\bar F(\sigma(x))$. A pattern $\sigma$
determines the set $\{x\in X:\sigma(x)=\sigma\}$, which is nonempty by the
definition of $\mathcal P_i$.

(iii) Each $\mathcal A_t$ is computed by testing at most $3^p$ patterns
against at most $2^p$ residual entries. A constraint problem whose relations
live on the bags of a tree decomposition is solved by nonserial dynamic
programming over $(T,(B_t))$ in the min-sum, max-sum and sum-product
semirings, with $O(p\,3^p)$ operations per bag \cite{Dechter1999}; each unary
cost is charged once, at the bag nearest the root that contains its variable.
The optimizer is unique if and only if exactly one pattern is admissible and
it has no $\ast$, by the disjointness in (ii). The dimension of $S$ is the
maximum number of $\ast$ entries on continuous coordinates. When $S$ is finite,
$|S|=\sum_\sigma\prod_{i:\sigma_i=\ast}|X_i\cap(l_i,u_i)|$. By (ii),
$\dist(y,S)^2=\min_\sigma\sum_ic_i(\sigma_i)$ over admissible $\sigma$, with
$c_i(\mathsf L)=(y_i-l_i)^2$, $c_i(\mathsf U)=(y_i-u_i)^2$ and
$c_i(\ast)=\dist(y_i,X_i)^2$. A linear objective attains its minimum over each
closed face at an endpoint pattern, so it is minimized by the same
recursion over $\{\mathsf L,\mathsf U\}$ labels. All numbers are rationals of
polynomial bit length.
\end{proof}

\begin{remark}[Two-label integer coordinates and multilinear objectives]\label{rem:endpointext}
If $X_i=\{l_i,l_i+1\}$, replacing $\frac12H_{ii}x_i^2$ by
$\frac12H_{ii}\bigl((2l_i+1)x_i-l_i(l_i+1)\bigr)$ does not change $F$ on $X$.
Hence the sign hypothesis is needed only for coordinates with at least three
feasible values; all binary quadratic programs are covered. The proofs also
apply verbatim to $F=\sum_af_a+\sum_i\psi_i(x_i)$ in which each $f_a$ is
affine in each of its variables separately and each $\psi_i$ is concave on
$[l_i,u_i]$. Then $\E f_a(Y)=f_a(x)$, and the diagonal term of
\eqref{eq:endpointid} becomes $\psi_i(x_i)$ minus its chord at $x_i$. That
difference vanishes at an interior point exactly when $\psi_i$ is affine on
$[l_i,u_i]$, which replaces the condition $H_{ii}<0$ in (a).
\end{remark}

Endpoint optimality for nonpositive diagonals is classical; Del Pia and
Khajavirad use it to separate binary-valued from genuinely continuous
variables \cite{DelPiaKhajavirad2026}. Min-sum messages and nonnegative tree
reparameterizations are also standard \cite{Dechter1999,WainwrightJaakkolaWillsky2005}.
The identity \eqref{eq:endpointid} records what these ingredients give for
the entire optimal set. Corollary~\ref{cor:facecsp} shows that the
decomposition also controls the geometry of $S$, not only the optimal value.

\subsubsection{The diagonal Lagrangian certificate}

From here to the end of the section all coordinates are continuous. A point
$\bar x\in X$ is a \emph{box KKT point} if $\partial_iF(\bar x)\ge0$ when
$\bar x_i=l_i$, $\partial_iF(\bar x)\le0$ when $\bar x_i=u_i$, and
$\partial_iF(\bar x)=0$ otherwise. For such a point put
\[
\lambda_i(\bar x)=\frac{2\,|\partial_iF(\bar x)|}{u_i-l_i},\qquad
\Lambda(\bar x)=\operatorname{diag}(\lambda_1(\bar x),\ldots,\lambda_n(\bar x)).
\]
For $\lambda\in\R^n_{\ge0}$ define the Lagrangian dual function
\[
\psi(\lambda)=\inf_{x\in\R^n}\Bigl[F(x)-\frac12\sum_i\lambda_i(x_i-l_i)(u_i-x_i)\Bigr].
\]

\begin{lemma}[Diagonal certificate, full optimal set and invariance]\label{lem:diagcert}
\begin{enumerate}
\item[(i)] For every box KKT point $\bar x$ and every $x\in\R^n$,
\begin{equation}\label{eq:diagid}
F(x)-F(\bar x)=\frac12(x-\bar x)^{\mathsf T}\bigl(H+\Lambda(\bar x)\bigr)(x-\bar x)
+\frac12\sum_i\lambda_i(\bar x)(x_i-l_i)(u_i-x_i).
\end{equation}
\item[(ii)] If $H+\Lambda(\bar x)\succeq0$, then $\bar x\in S$ and
\begin{equation}\label{eq:diagset}
S=\bigl\{x\in X:\ (H+\Lambda(\bar x))(x-\bar x)=0,\ \
\lambda_i(\bar x)(x_i-l_i)(u_i-x_i)=0\ \text{for all }i\bigr\}.
\end{equation}
\item[(iii)] $\psi(\lambda)\le\OPT$ for every $\lambda\ge0$, and the
following are equivalent: (a) $H+\Lambda(\bar x)\succeq0$ for some
$\bar x\in S$; (b) $H+\Lambda(\bar x)\succeq0$ for every $\bar x\in S$;
(c) $\psi(\lambda)=\OPT$ for some $\lambda\ge0$. When they hold, the vector
$\lambda(\bar x)$ is the same for every $\bar x\in S$, and it is the unique
maximizer of $\psi$.
\end{enumerate}
\end{lemma}

\begin{proof}
(i) Both sides are quadratic polynomials in $x$ that vanish at $\bar x$.
Their Hessians are $H$ and $H+\Lambda-\Lambda=H$. The $i$th partial
derivative of the right side at $\bar x$ is
$\frac12\lambda_i(\bar x)(u_i+l_i-2\bar x_i)$. This equals
$|\partial_iF(\bar x)|=\partial_iF(\bar x)$ when $\bar x_i=l_i$, equals
$-|\partial_iF(\bar x)|=\partial_iF(\bar x)$ when $\bar x_i=u_i$, and is
$0=\partial_iF(\bar x)$ at an interior coordinate. Two quadratics with equal
value, gradient and Hessian at one point coincide.

(ii) On $X$ both terms on the right of \eqref{eq:diagid} are nonnegative,
so $\bar x$ is optimal, and $x\in X$ is optimal exactly when both vanish. A
positive semidefinite quadratic form vanishes exactly on the kernel of its
matrix.

(iii) For $x\in X$ the bracket in $\psi$ is at most $F(x)$, so
$\psi(\lambda)\le\OPT$. If (a) holds at $\bar x$, then by (i) the bracket with
$\lambda=\lambda(\bar x)$ equals $F(\bar x)+\frac12(x-\bar x)^{\mathsf T}(H+\Lambda)(x-\bar x)$,
whose infimum over $\R^n$ is $F(\bar x)=\OPT$; this gives (c). Suppose (c)
holds for some $\lambda$, and let $\bar x\in S$ be arbitrary. Write
$\mathcal L(x)$ for the bracket. Then
$\OPT=\psi(\lambda)\le\mathcal L(\bar x)\le F(\bar x)=\OPT$. Hence
$\lambda_i(\bar x_i-l_i)(u_i-\bar x_i)=0$ for all $i$, and $\bar x$ minimizes
the quadratic $\mathcal L$ over $\R^n$. So
$\nabla^2\mathcal L=H+\operatorname{diag}(\lambda)\succeq0$ and
$0=\partial_i\mathcal L(\bar x)=\partial_iF(\bar x)+\lambda_i\bigl(\bar x_i-\tfrac{l_i+u_i}2\bigr)$.
At an interior coordinate complementarity gives $\lambda_i=0$. At
$\bar x_i=l_i$ the stationarity equation gives
$\lambda_i=2\partial_iF(\bar x)/(u_i-l_i)=\lambda_i(\bar x)$, and similarly at
$\bar x_i=u_i$. Thus $\lambda=\lambda(\bar x)$ and
$H+\Lambda(\bar x)\succeq0$. Since $\bar x\in S$ was arbitrary, this proves (b)
and the uniqueness and invariance statements. Finally (b) implies (a)
because $S\ne\emptyset$.
\end{proof}

\begin{remark}\label{rem:shor}
Call the instances satisfying (a)--(c) the \emph{diagonal certificate class}
$\mathcal C$. Convex box QPs belong to it, since $\Lambda(\bar x)\succeq0$. It also contains indefinite instances with tilted and
disconnected optimal sets. For example $F=(x-y+t/2)^2+\frac18t(1-t)$ on
$[0,1]^3$ has diagonal $(2,2,\frac14)$, the two optimal segments
$\{(v,v,0)\}$ and $\{(v,v+\frac12,1):0\le v\le\frac12\}$, the multiplier
$\Lambda=\operatorname{diag}(0,0,\frac14)$ at every optimizer, and
$H+\Lambda=2(1,-1,\tfrac12)^{\mathsf T}(1,-1,\tfrac12)\succeq0$. The function
$\psi$ is the Lagrangian dual for the constraints
$(x_i-l_i)(u_i-x_i)\ge0$, so condition (c) says that this dual, equivalently
the basic Shor semidefinite relaxation, is exact. Sufficient global
optimality conditions of this box-multiplier form are classical
\cite{JeyakumarRubinovWu2006,LiWuQuan2015}, and stronger relaxations with
complete exactness characterizations are studied in \cite{QiuYildirim2023}.
Lemma~\ref{lem:diagcert} is therefore not a new certificate. What we use is
its invariance: whichever optimizer a search finds, the same rational test
accepts it.
\end{remark}

\subsubsection{Function values cannot replace a certificate under set growth}

\begin{proposition}[Value-oracle barrier for unknown set growth]\label{prop:oraclebarrier}
Fix $k\ge1$, $X=[0,1]^k$ and the coordinate-curvature bound $L=1$. Consider a
deterministic algorithm that queries exact values, gradients and Hessians of
an unknown $C^2$ function $f$ at adaptively chosen points and outputs either a
rational interval of width at most $\varepsilon$ containing $\min_Xf$, or a
point $\hat x$ with $f(\hat x)-\min_Xf\le\varepsilon$. Suppose it is correct
for every $C^2$ function with $\partial_{ii}f\le1$ on $X$ for all $i$ that
satisfies \eqref{eq:setgrowth} for some $g>0$. Then on $f\equiv0$, for which
$\kappa=1$, it makes at least $\bigl(1/(2\sqrt{6\varepsilon})-1\bigr)^k-1$
queries.
\end{proposition}

\begin{proof}
For $c\in X$ and $0<w\le1$ let
$f_{c,w}(x)=-\frac{w^2}6\bigl(1-|x-c|^2/w^2\bigr)^3$ when $|x-c|\le w$ and
$f_{c,w}(x)=0$ otherwise. In the variable $z=(x-c)/w$ with $\rho=|z|^2$,
\[
\partial_if_{c,w}=w\,z_i(1-\rho)^2,\qquad
\partial_{ij}f_{c,w}=\delta_{ij}(1-\rho)^2-4z_iz_j(1-\rho).
\]
Value, gradient and Hessian vanish on $\rho=1$, so $f_{c,w}$ is $C^2$, and
$\partial_{ii}f_{c,w}=(1-\rho)(1-\rho-4z_i^2)\le1$. The minimizer $c$ is
unique with value $-w^2/6$. Inside the ball,
$f_{c,w}+\frac{w^2}6=\frac{w^2}6\bigl(1-(1-\rho)^3\bigr)\ge\frac{w^2}6\rho
=\frac{|x-c|^2}6$, because $1-(1-\rho)^3-\rho=\rho(1-\rho)(2-\rho)\ge0$.
Outside it the gap is $w^2/6\ge\frac{w^2}{6k}|x-c|^2$. Hence
\eqref{eq:setgrowth} holds with $g=w^2/(6k)$. The zero function satisfies it
with $g=1$ and $S=X$.

Run the algorithm on $f\equiv0$; let $q$ be its number of queries and add
the output point, if any, to the query points. Suppose
$q<\bigl(1/(2\sqrt{6\varepsilon})-1\bigr)^k-1$, and let
$M=\lfloor(q+1)^{1/k}\rfloor+1$. Then $M^k>q+1$ and $M<1/(2\sqrt{6\varepsilon})$.
Divide $X$ into $M^k$ closed subcubes of side $1/M$. The interior of one of
them contains none of the at most $q+1$ points. Let $c$ be its center and
choose $w$ with $\sqrt{6\varepsilon}<w<1/(2M)$. The closed ball of radius $w$
lies in the open subcube, so every oracle answer for $f_{c,w}$ at the
recorded points equals the answer for $0$. By determinism the algorithm
produces the same transcript and output on $f_{c,w}$. An interval of width at
most $\varepsilon$ containing $0$ does not contain $-w^2/6<-\varepsilon$. An
output point outside the ball has gap $w^2/6>\varepsilon$ for $f_{c,w}$.
Either way the algorithm is wrong on an admissible function.
\end{proof}

The filtered-grid algorithm of Section~\ref{sec:rate} uses only exact values
and the curvature bound, yet needs no accuracy-dependent number of states under
point growth. Proposition~\ref{prop:oraclebarrier} shows that no algorithm in
that information model can do the same under growth toward an unknown set,
even with $k=p=1$ and $\kappa=1$. The barrier concerns information, not the
input class: the alternatives are not quadratics. It is the classical
hidden-bump argument of information-based complexity, adapted to set growth.
An algorithm for explicit quadratics must therefore exploit the coefficients,
for example through a certificate such as Lemma~\ref{lem:diagcert}.

\subsubsection{A proximal grid under a trusted growth guess}

The following generator finds a point near $S$ when a guessed ratio
$K\ge1$ satisfies $K\ge L/g$. Its work is bounded for every guess, but its
conclusion is not checkable, so it is used only to propose candidates.

Given a rational $K\ge1$, let $\mu\ge2$ be the least integer with
$4^{-\mu}K\le1/8$, and put
\[
\theta=2^{-\mu},\qquad\eta=\frac{L\theta^2}4,\qquad
\varrho=2\bigl\lceil\sqrt{Kn}\,\bigr\rceil,\qquad
K_\theta=100\,\theta^{-1}\lceil\log_2(n+2)\rceil .
\]
Then $\theta\le1/4$, $\theta^2K\le1/8$ and $\theta^{-1}\le6\sqrt K$.
A \emph{proximal stage} with center $c\in X$ and mesh $h>0$ does the
following. Form the fresh box $B=X\cap\prod_i[c_i-\varrho h,c_i+\varrho h]$,
rounding integer endpoints inward. Build in each $B_i$ the outward geometric
grid $G_i$ from $c_i$ of Section~\ref{sec:rate}, with step $h+\theta t$
(continuous) or $\max\{1,\lfloor h+\theta t\rfloor\}$ (integer) at distance
$t$, clipped at the endpoints of $B_i$. With the corrections $d_i$ and $D$ of
\eqref{eq:correction} on these grids, minimize
\[
P(y)=F(y)-D(y)+\eta\norm{y-c}^2\qquad\Bigl(y\in\prod_iG_i\Bigr)
\]
by the junction-tree recursion \eqref{eq:messages}; correction and proximal
terms are unary, so the decomposition is unchanged. The output is a
minimizer $y^+$. The \emph{proximal iteration with guess $K$} starts at the
lower endpoint vector $c^0=l$, uses $h_j=s2^{-j}$, outputs $y^j$ at stage $j$,
and sets $c^{j+1}=y^j$. Every center is feasible, and every box is built from
$X$ rather than from the previous box.

\begin{lemma}[Proximal stage]\label{lem:proximal}
\begin{enumerate}
\item[(i)] For every guess, every grid has at most $K_\theta$ nodes.
\item[(ii)] Let $g$ satisfy \eqref{eq:setgrowth} with $L/g\le K$. If
$\dist(c,S)^2\le4Knh^2$, then
\[
F(y^+)-\OPT\le\tfrac{99}{256}Lnh^2,\qquad
\dist(y^+,S)^2\le\tfrac{99}{256}Knh^2 .
\]
\item[(iii)] Under the hypothesis of (ii) on $g$, the proximal iteration
satisfies $\dist(y^j,S)^2\le\frac{99}{256}Knh_j^2$ for every $j\ge0$.
\item[(iv)] Stage $j$ costs $O\bigl(p(N+|\mathcal A|)K_\theta^p\bigr)$ table
operations on rationals of bit length $\poly(I+j+\mu K_\theta)$, for every
guess.
\end{enumerate}
\end{lemma}

\begin{proof}
(i) Each endpoint of $B_i$ is within $\varrho h$ of $c_i$. As in the proof of
Lemma~\ref{lem:states}, one side of a continuous grid has at most
$1+\lceil\ln(1+\theta\varrho)/\ln(1+\theta)\rceil$ nodes, and one side of an
integer grid at most $1+\lceil\ln(1+\theta\varrho h/H)/\ln(1+\theta/3)\rceil$
with $H=\max\{h,1\}\ge h$. Now
$\theta\varrho\le2\theta\sqrt{Kn}+2\theta\le\sqrt{n/2}+\frac12$, and
$\frac32+\sqrt{n/2}\le n+2$. Using $\ln(1+\theta)\ge\ln(1+\theta/3)\ge\theta/4$
and $\ln(n+2)\le\log_2(n+2)$, each side has at most
$2+4\theta^{-1}\log_2(n+2)$ nodes. With the center, $|G_i|\le5+8\theta^{-1}
\lceil\log_2(n+2)\rceil\le10\,\theta^{-1}\lceil\log_2(n+2)\rceil\le K_\theta$,
since $\theta^{-1}\ge4$.

(ii) Let $s^*\in S$ be nearest to $c$ (it exists since $S$ is compact) and
$E=\norm{s^*-c}^2\le4Knh^2$. Each $|s^*_i-c_i|\le2\sqrt{Kn}\,h\le\varrho h$,
and integer coordinates of $s^*$ are integers, so $s^*\in B$. Let $Y$ be the
independent mean-preserving rounding of $s^*$ to its enclosing grid
intervals. Lemma~\ref{lem:round} on the box $B$ gives
$\E[F(Y)-D(Y)]\le F(s^*)=\OPT$. Independence gives
$\E\norm{Y-c}^2=E+\sum_i\operatorname{Var}Y_i$. If $s^*_i$ is not a node, its
enclosing interval was generated from its inner endpoint, at distance at most
$|s^*_i-c_i|$ from $c_i$, so its length is at most $h+\theta|s^*_i-c_i|$;
an integer coordinate inside a unit interval is impossible, and clipping only
shortens. Hence
$\operatorname{Var}Y_i\le\frac14(h+\theta|s^*_i-c_i|)^2\le\frac12(h^2+\theta^2|s^*_i-c_i|^2)$
and $\sum_i\operatorname{Var}Y_i\le\frac12(nh^2+\theta^2E)$. Since a minimum
is at most an expectation,
\[
P(y^+)\le\OPT+\eta\Bigl[\bigl(1+\tfrac{\theta^2}2\bigr)E+\tfrac{nh^2}2\Bigr].
\]
By \eqref{eq:mesh}, $w_i(v)\le h+\theta|v-c_i|$ for every node, so
$d_i(v)\le\frac L4h^2+\eta|v-c_i|^2$ and
$D(y^+)\le\frac L4nh^2+\eta\norm{y^+-c}^2$. Therefore
\begin{align*}
F(y^+)-\OPT&=P(y^+)-\OPT+D(y^+)-\eta\norm{y^+-c}^2\\
&\le\frac{Lnh^2}4\Bigl[1+4K\theta^2\bigl(1+\tfrac{\theta^2}2\bigr)+\tfrac{\theta^2}2\Bigr]
\le\frac{Lnh^2}4\cdot\frac{99}{64},
\end{align*}
using $E\le4Knh^2$, $4K\theta^2\le\frac12$ and $\theta^2\le\frac1{16}$.
Growth gives $\dist(y^+,S)^2\le(F(y^+)-\OPT)/g\le\frac{99}{256}(L/g)nh^2\le
\frac{99}{256}Knh^2$.

(iii) Initially $\dist(l,S)^2\le ns^2\le4Knh_0^2$. If
$\dist(y^j,S)^2\le\frac{99}{256}Knh_j^2$, then
$\dist(c^{j+1},S)^2\le Knh_j^2/2=2Knh_{j+1}^2\le4Knh_{j+1}^2$. Apply (ii).

(iv) Let $\omega_0$ be a common denominator of $H,b,c,L$ and the endpoints.
By induction, every stage-$j$ grid coordinate has a denominator dividing
$\omega_02^{j+\mu K_\theta}$: the center comes from stage $j-1$; the new
endpoints $c_i\pm\varrho h_j$ have denominators dividing $\omega_02^j$; an
unclipped continuous offset is
$h_j\bigl[(2^\mu+1)^k-2^{\mu k}\bigr]/2^{\mu(k-1)}$ with $k\le K_\theta$;
integer steps are integral; clipping selects an endpoint of $B_i$. Values of
$F$, corrections and the proximal term then have a common denominator
dividing $8\omega_0^22^{2\mu+2}\bigl(\omega_02^{j+\mu K_\theta}\bigr)^2$, and
bounded magnitudes because all points stay in $\hat X$. Messages are sums of
such values. The table count is \eqref{eq:dpwork} with $K_\theta$ labels.
\end{proof}

\begin{remark}[A false guess]\label{rem:falseguess}
Without $L/g\le K$ the iteration may converge to a nonglobal point. For
$F=x^2+y^2-3xy+\frac{63}{128}(x+y)$ on $[0,1]^2$, put $t=x+y$; then
$F\ge t^2-\frac54t^2+\frac{63}{128}t$, a concave function of $t\in[0,2]$, so
$(1,1)$ is the unique optimizer with $\OPT=-\frac1{64}$. The origin is a box
KKT point with $F=0$. Exact computation shows that the iteration with
$K\in\{1,2,4\}$ returns the origin at every stage $j\le33$, which covers the
budgets $J_K$ of Theorem~\ref{thm:diagdiscovery}. Lemma~\ref{lem:diagcert}
rejects it: $H+\Lambda=\bigl(\begin{smallmatrix}191/64&-3\\-3&191/64\end{smallmatrix}\bigr)$
has eigenvalue $-\frac1{64}$. At $(1,1)$ the matrix has diagonal $\frac{193}{64}$
and is positive definite.
\end{remark}

\subsubsection{Stationary-face recovery}

This lemma is the recovery step inside the proof of
Theorem~\ref{thm:generalfinite}, stated separately because it is used twice.
Let $\omega$ be the least positive integer for which $\omega H$, $\omega b$
and $\omega l,\omega u$ are integral, and put
\[
C_H=\max\{1,\max_{ij}|\omega H_{ij}|\},\qquad\Delta=(nC_H)^n,\qquad
R=\omega\Delta,\qquad\tau=\frac1{4nR}.
\]
These numbers have bit length polynomial in $I$. For $y\in X$, the
\emph{selected face} fixes every integer coordinate to $y_i$, fixes a
continuous coordinate to $l_i$ if $y_i-l_i\le\tau$ and to $u_i$ if
$u_i-y_i\le\tau$, and leaves the other continuous coordinates free; let $J$ be
the free set and
\[
\Pi(y)=\{x\in\hat X:\ \text{selected coordinates fixed},\ \ \partial_jF(x)=0\ (j\in J)\}.
\]

\begin{lemma}[Stationary-face recovery]\label{lem:recovery}
Every vertex of a polytope defined by fixing coordinates to rational values
of denominator dividing $\omega$, the bounds of $\hat X$, and stationarity
equations $\partial_jF(x)=0$ has a common coordinate denominator at most $R$.
If $y\in X$ and $\dist(y,S)\le\tau/2$, then $\Pi(y)$ is nonempty and every
point of $\Pi(y)$ is a global optimizer.
\end{lemma}

\begin{proof}
In $z=\omega x$, a fixing or bound row is a unit row with integral right side,
and a stationarity row reads $(\omega H)_{j\cdot}z=-\omega^2b_j$, with integral
entries of magnitude at most $C_H$. A vertex solves $n$ linearly independent
such rows $Az=\beta$. Expanding $\det A$ along its unit rows leaves a
$k\times k$ minor of $\omega H$ with $k\le n$, of magnitude at most
$k!\,C_H^k\le\Delta$. By Cramer's rule $z$ has common denominator $|\det A|$,
so $x=z/\omega$ has one at most $\omega\Delta=R$.

Let $s^*\in S$ be nearest to $y$. Since $\norm{y-s^*}\le\tau/2<1$, the integer
coordinates agree. Every positive width is at least $1/\omega>2\tau$, so no
coordinate is selected at both bounds. Let $A^*$ be the continuous coordinates
at a bound in $s^*$ and $J_0$ the others. Each $i\in A^*$ has $y_i$ within
$\tau/2$ of the same bound, so it is selected there; hence $J\subseteq J_0$.
Let
\[
P=\{x\in\hat X:\ x_i=s^*_i\ (i\text{ integer or }i\in A^*),\ \ \partial_jF(x)=0\ (j\in J_0)\}.
\]
Optimality of $s^*$ with two-sided feasible moves in $J_0$ gives $s^*\in P$.
For $x\in P$ the difference $d=x-s^*$ is supported on $J_0$ and
$H_{J_0J_0}d_{J_0}=\nabla_{J_0}F(x)-\nabla_{J_0}F(s^*)=0$, so
$F(x)-F(s^*)=\nabla F(s^*)^{\mathsf T}d+\frac12d^{\mathsf T}Hd=0$: every point
of $P$ is optimal. Let $q(x)$ be the sum of the slacks $x_i-l_i$ or $u_i-x_i$
of the coordinates in $J_0\setminus J$ to their selected bounds. At $s^*$ each
such slack is at most $\tau+\tau/2$, so $q(s^*)\le\frac32n\tau=\frac3{8R}$. The
linear function $q$ attains its minimum over the nonempty polytope $P$ at a
vertex. There $\omega q$ is an integral combination of the coordinates of
$z=\omega x$, which have a common denominator at most $\Delta$, plus an
integer; so a positive minimum would be at least $1/(\omega\Delta)=1/R$. Hence some $s'\in P$ has $q(s')=0$; it lies in $\Pi(y)$ because
$J\subseteq J_0$. For any $x\in\Pi(y)$, $d=x-s'$ is supported on $J$ and
$H_{JJ}d_J=0$, so $F(x)=F(s')+\nabla_JF(s')^{\mathsf T}d_J+\frac12d_J^{\mathsf T}H_{JJ}d_J=\OPT$.
\end{proof}

\subsubsection{Exact discovery for the certificate class with unknown growth}

\begin{theorem}[Unknown-growth discovery of a certified optimal set]\label{thm:diagdiscovery}
Let $X$ be a continuous box. Consider the following procedure. For
$K=1,2,4,\ldots$: let $J_K$ be the least $j\ge0$ with $4^j\tau^2\ge4Kns^2$;
run the proximal iteration with guess $K$ for stages $0,\ldots,J_K$, and let
$y=y^{J_K}$. If $\Pi(y)$ is nonempty, compute a vertex $\bar x$ of it by
exact linear programming. If $\bar x$ is a box KKT point and
$H+\Lambda(\bar x)\succeq0$, output $\bar x$, $\OPT=F(\bar x)$, $\Lambda(\bar x)$
and the description \eqref{eq:diagset}; otherwise continue with $2K$.
\begin{enumerate}
\item[(a)] Every output is correct, with no assumption on the instance.
\item[(b)] If $F\in\mathcal C$ and $g>0$ is any set-growth constant, the
procedure stops no later than the first guess $K\ge L/g$, so with $K\le2\kappa$.
\item[(c)] Its running time is $f(p,\kappa)\poly(I)$ bit operations with an
absolute polynomial exponent, and the output has bit length $\poly(I)$.
\end{enumerate}
Neither $S$, nor $g$, nor a bound on $\kappa$ is supplied.
\end{theorem}

\begin{proof}
(a) Acceptance requires exactly the hypotheses of
Lemma~\ref{lem:diagcert}(ii), checked in rational arithmetic; the positive
semidefiniteness test is symmetric Gaussian elimination that rejects a
negative pivot and a zero pivot with a nonzero off-diagonal entry in its row.

(b) If $K\ge L/g$, Lemma~\ref{lem:proximal}(iii) gives
$\dist(y,S)^2\le\frac{99}{256}Kns^24^{-J_K}<\tau^2/4$. By
Lemma~\ref{lem:recovery}, $\Pi(y)$ is nonempty and its vertex $\bar x$ is
optimal, hence a box KKT point. By Lemma~\ref{lem:diagcert}(iii), (a) implies
(b), so $H+\Lambda(\bar x)\succeq0$ and the test accepts, whichever optimal
component $\bar x$ belongs to.

(c) At most $1+\log_2(2\kappa)$ guesses run. For each,
$J_K+1=O(\log(Kns^2/\tau^2))=\poly(I)+O(\log K)$, and by
Lemma~\ref{lem:proximal}(i),(iv) a stage costs
$O(p(N+|\mathcal A|)K_\theta^p)$ operations on numbers of bit length
$\poly(I+J_K+\mu K_\theta)$, where $\theta^{-1}\le6\sqrt K$ and $\mu=O(\log K)$.
As in Theorem~\ref{thm:approx}, \eqref{eq:logabsorb} absorbs
$\lceil\log_2(n+2)\rceil^p$ into $(C_0p)^p(n+2)$, and every factor depending
only on $p$ and $K\le2\kappa$ enters $f$. Selection uses original data, so
$\Pi(y)$ has input-size description; exact linear programming returns a vertex
in polynomial time, and the vertex has denominators at most $R$ by
Lemma~\ref{lem:recovery}. The KKT and semidefiniteness tests are polynomial.
\end{proof}

On $\mathcal C$ the procedure always terminates, because a set-growth
constant always exists \cite[Theorem~3.3]{LuoSturm2000}; $\kappa$ governs only
the running time. Outside $\mathcal C$ it may run forever, so an
implementation needs a work limit, and an unsuccessful guess certifies
nothing about membership. For fixed $p$ the bound is polynomial in $\kappa$
and $I$: the table size is $O\bigl((600\sqrt K)^p(C_0p)^p(n+2)\bigr)$, and the
bit lengths are polynomial in $I$ and $\sqrt K\log K$.

\begin{proposition}[Discovery within $\mathcal C$ is hard without the growth parameter]\label{prop:sshard}
For integers $a_1,\ldots,a_m$ and $T$ with $|T|\le A=\sum_i|a_i|$, let
$\sigma_0=0$ and
\[
F(x,\sigma)=\sum_{i=1}^m(\sigma_i-\sigma_{i-1}-a_ix_i)^2+(\sigma_m-T)^2
+\sum_{i=1}^mx_i(1-x_i)
\]
on $x\in[0,1]^m$, $\sigma\in[-A,A]^m$. The bags $\{\sigma_1,x_1\}$ and
$\{\sigma_{i-1},\sigma_i,x_i\}$ $(2\le i\le m)$ form a path decomposition of
size at most $3$. If some $\hat x\in\{0,1\}^m$ has $\sum_ia_i\hat x_i=T$, then
$F\in\mathcal C$, $\OPT=0$, and $S$ consists exactly of the points
$(\hat x,\hat\sigma)$ with $\hat x$ such a solution and $\hat\sigma$ its
partial sums. Otherwise $\OPT>0$. Consequently, unless $\mathrm P=\mathrm{NP}$,
no polynomial-time algorithm returns an exact optimizer for every instance of
$\mathcal C$ with bag size three, even though $\OPT=0$ is known in advance for
these instances, and the ratio $\kappa$ of these instances is not bounded by
any polynomial in their input length.
\end{proposition}

\begin{proof}
$F\ge0$ on the box, and $F=0$ exactly when all squares and all products
$x_i(1-x_i)$ vanish, that is, for binary $x$ with partial sums $\sigma$ and
$\sigma_m=T$; the partial sums lie in $[-A,A]$. If no solution exists, $F>0$
on the compact box. In the feasible case let $\hat z=(\hat x,\hat\sigma)$ be
optimal. The sum of
squares equals $\frac12(z-\hat z)^{\mathsf T}M(z-\hat z)$ with $M\succeq0$,
because each square is of an affine function vanishing at $\hat z$. Also
$\sum_ix_i(1-x_i)=\frac12\sum_i2(x_i-0)(1-x_i)$. So
$F(z)-F(\hat z)=\frac12(z-\hat z)^{\mathsf T}M(z-\hat z)+\frac12\sum_i2x_i(1-x_i)$.
At $\hat z$ the gradient is $1-2\hat x_i=\pm1$ in $x_i$ and $0$ in
$\sigma$, so $\hat z$ is a box KKT point with $\lambda_i(\hat z)=2$ on $x$ and
$0$ on $\sigma$, and $H+\Lambda(\hat z)=M\succeq0$. Thus $F\in\mathcal C$.
Subset Sum with binary-encoded integers is NP-complete. An algorithm as in the
statement, stopped at its polynomial time bound and followed by a check of its
output, would decide it. By the remark after Theorem~\ref{thm:diagdiscovery},
the same would follow if $\kappa$ were polynomially bounded on these instances.
\end{proof}

The optimal value of an instance of $\mathcal C$ equals $\max\psi$, the value
of a convex relaxation. Proposition~\ref{prop:sshard} shows that the
difficulty is in locating an optimizer, and that a parameter such as $\kappa$
is needed for an efficient exact search. Corollary~\ref{cor:facecsp} has no
counterpart here. With $T=0$ and integers $a_i$ of both signs, the origin is an
optimizer of an instance of $\mathcal C$, and it is the unique optimizer
exactly when no nonempty subset of the $a_i$ sums to zero, which is an
NP-complete question. So even with one optimizer and the full description
\eqref{eq:diagset} in hand, deciding uniqueness is hard.

% ---------------------------------------------------------------------------
\section{Exact implicit output for polynomial boundary optima}\label{app:boundary}

This appendix makes precise what exact output means for the explicit
polynomials of Section~\ref{sec:polynomial} when the optimizer can be
irrational and lie on the boundary. Throughout, $F$ is an explicit rational
polynomial of fixed degree $d\ge2$ with factor scopes in the bags, on a mixed
box $X$, with a verified curvature bound $L$. Let $\mathcal C_{\rm c}$ be the
continuous coordinates. Fix rationals $C_2,C_3\ge1$ with
\[
\sum_{j\in\mathcal C_{\rm c}}\sup_{\hat X}|\partial_{ij}F|\le C_2,\qquad
\sum_{j,k\in\mathcal C_{\rm c}}\sup_{\hat X}|\partial_{ijk}F|\le C_3
\qquad(i\in\mathcal C_{\rm c}).
\]
Monomialwise interval bounds give such numbers with bit length $\poly(I)$
at fixed $d$.

\paragraph{Three sound tests.}
The pruned-grid trials of Section~\ref{sec:rate} are run with one extra filter.

\begin{lemma}[Integer-label filter]\label{lem:labelfilter}
(i) Suppose at some stage every grid interval of an integer coordinate $i$
has length one. Then deleting the labels $v$ with $m_i(v)>U$ and taking the
hull of the remaining labels removes no point with objective at most $U$, and
keeps the incumbent and the corrected-grid minimizer; the filtering-history
argument of Proposition~\ref{prop:filter} remains valid.
(ii) Assume a unique optimizer $x^*$, point growth \eqref{eq:growth}, and
$\theta^2\le1/(8\kappa)$. At the first stage with
$h_j\le1/(10\sqrt{n\kappa})$, every integer coordinate is fixed to $x^*_i$.
\end{lemma}

\begin{proof}
(i) Every interval adjacent to $v$ is a unit interval, so $d_i(v)=0$. For
feasible $x$ with $x_i=v$, round the other coordinates as in
Lemma~\ref{lem:round}; coordinate $i$ stays at $v$. Then
$F(x)\ge\E Q(Y)\ge m_i(v)$. Hence a deleted label carries no point with
value at most $U$, in particular no optimizer and not the incumbent. The
corrected-grid minimizer $y$ survives because $m_i(y_i)\le Q(y)=\beta\le\OPT\le U$.

(ii) If $j=0$ then $s<1$ and no integer coordinate is left. Otherwise, by
Lemma~\ref{lem:localize} at stage $j-1$, the current integer domain lies
within $1+10\sqrt{n\kappa}\,h_j\le2$ of the center. All steps
$\max\{1,\lfloor h_j+\theta t\rfloor\}$ with $t\le2$ equal one, so (i) applies.
A retained label $v$ has a grid assignment $z$ with $z_i=v$ and
$Q(z)\le U_j$; the derivation of \eqref{eq:witnessradius} applies to it and
gives $\norm{z-x^*}^2\le\frac{22}{15}\kappa nh_j^2<1$, so $v=x^*_i$.
The extra filter keeps $x^*$ and $y_j$, so Lemmas~\ref{lem:contraction}
and~\ref{lem:localize} still apply.
\end{proof}

\begin{lemma}[Monotone face reduction]\label{lem:monotone}
Let $B$ be a product box with fixed integer coordinates, $i$ a free
continuous coordinate whose original bound $l_i$ is an endpoint of $B_i$, and
suppose $\partial_iF\ge0$ on $B$. Then $\min_BF=\min_{B\cap\{x_i=l_i\}}F$, and
if $\partial_iF>0$ on $B$, every minimizer over $B$ has $x_i=l_i$. The same
holds for $u_i$ with $\partial_iF\le0$. With $c$ the midpoint and $r$ the
largest half-width of the free continuous coordinates of $B$,
$\partial_iF(c)-C_2r>0$ certifies $\partial_iF>0$ on $B$, and
$\partial_iF(c)+C_2r<0$ certifies $\partial_iF<0$.
\end{lemma}

\begin{proof}
Replacing $x_i$ by $l_i$ keeps a point of $B$ in $B$ and, by the mean value
theorem along the segment, does not increase $F$, and decreases it under a
strict sign unless $x_i=l_i$. For the test,
$|\partial_iF(x)-\partial_iF(c)|\le\sum_j\sup|\partial_{ij}F|\,|x_j-c_j|\le C_2r$.
\end{proof}

Any other sound sign enclosure may be used as well, for example Bernstein
coefficient bounds of each factor derivative on its own scope, which have at
most $d^p$ coefficients. Such monotonicity tests are standard
\cite{ArayaTrombettoniNeveu2010}; the termination proof below uses only the
midpoint test.

\begin{lemma}[Uniform convexity patch]\label{lem:patch}
Let $B$ be a box with fixed integer coordinates and some fixed continuous
coordinates, $J$ its free coordinates, $c$ its midpoint and $r$ the largest
free half-width. If $H_{JJ}(c)-(C_3r+\sigma)\mathrm{Id}\succ0$ for some
$\sigma>0$, then $F$ restricted to $B$ is strongly convex in $x_J$ and has a
unique minimizer over $B$.
\end{lemma}

\begin{proof}
Each entry satisfies $|H_{jk}(x)-H_{jk}(c)|\le r\sum_l\sup|\partial_{jkl}F|$,
so every row sum of $|H_{JJ}(x)-H_{JJ}(c)|$ is at most $C_3r$. The spectral
norm of a symmetric matrix is at most its largest absolute row sum. Hence
$H_{JJ}(x)\succ\sigma\mathrm{Id}$ on $B$.
\end{proof}

\paragraph{The procedure.}
Trial $\mu\ge2$ runs the capped pruned-grid algorithm with $\theta=2^{-\mu}$,
the cap \eqref{eq:statecap}, and the filter of Lemma~\ref{lem:labelfilter},
for stages $j=0,1,2,\ldots$ without a stage limit, aborting if the cap would be
exceeded. After a stage it succeeds if the certified gap is zero (the
incumbent is then an exact rational optimizer). Otherwise, if every integer
coordinate of the retained box is a singleton, it copies the retained box,
fixes the integer coordinates, applies sound reductions of
Lemma~\ref{lem:monotone} until none applies, and succeeds if no free
coordinate remains or if the test of Lemma~\ref{lem:patch} holds with
$\sigma_\mu=L/4^{\mu+1}$. The copy never influences later pruning. The
procedure runs phases $q=2,3,\ldots$; in phase $q$ it restarts trial $\mu$ for
each $2\le\mu\le q$ and stops it after $2^{q-\mu}$ bit operations, stopping at
the first success.

\begin{theorem}[Exact implicit boundary output with a complementarity margin]\label{thm:boundary}
(a) On success, the output is a fixed integer assignment, a set of continuous
coordinates fixed at original bounds, and a rational box for the remaining
coordinates on which $F$ is strongly convex; the unique minimizer $\hat x$ of
$F$ over this face box is a global optimizer of the original problem. If no
coordinate is free, $\hat x$ is an explicit rational optimizer. This holds
with no assumption on the instance.

(b) Suppose $F$ has a unique optimizer $x^*$ with point growth
\eqref{eq:growth}, and every continuous coordinate with $x^*_i\in\{l_i,u_i\}$
is strictly complementary: its inward derivative, $\partial_iF(x^*)$ at a
lower bound or $-\partial_iF(x^*)$ at an upper bound, is positive. Let
$\gamma$ be the smallest such inward derivative and
$B_\gamma=\lceil\log_2\max\{2,C_2/\gamma\}\rceil$ ($B_\gamma=1$ if there is
no active continuous coordinate). Then the procedure succeeds after
$f_d(p,\kappa)\poly(I+B_\gamma)$ bit operations. Neither $x^*$, $g$, nor
$\gamma$ is supplied.
\end{theorem}

\begin{proof}
(a) The retained box contains every optimizer
(Proposition~\ref{prop:filter} and Lemma~\ref{lem:labelfilter}), so after
fixing its singleton integer coordinates its minimum is $\OPT$. Each
reduction preserves the minimum over the current box
(Lemma~\ref{lem:monotone}). On the final face box $F$ has minimum $\OPT$ and,
by Lemma~\ref{lem:patch}, a unique minimizer, which is therefore globally
optimal. A zero certified gap makes the incumbent optimal.

(b) Let $\mu_*$ be the least $\mu\ge2$ with $4^\mu\ge8\kappa$, so that
$4^{\mu_*}\le32\kappa$ and trial $\mu_*$ satisfies the hypotheses of
Lemmas~\ref{lem:contraction}--\ref{lem:states}; it never reaches its cap. Let
$j_*$ be the least stage with
\[
h_j\le\frac1{10\sqrt{n\kappa}},\qquad
5\sqrt{n\kappa}\,h_j\le\frac\gamma{4C_2},\qquad
5\sqrt{n\kappa}\,h_j\le\frac L{8\cdot4^{\mu_*}C_3},
\]
omitting the first condition without integer coordinates and the second
without active continuous coordinates. Then
$j_*=\poly(I)+O(\log\kappa+B_\gamma)$, since $\log s$, $\log C_2$, $\log C_3$
and $\log(1/L)$ are polynomial in $I$. We show that trial $\mu_*$ succeeds by
stage $j_*$.

At stage $j_*$ the integer coordinates are fixed to $x^*$
(Lemma~\ref{lem:labelfilter}) and every continuous retained interval has
half-width $r\le5\sqrt{n\kappa}\,h_{j_*}$ (Lemma~\ref{lem:localize}). A sound
reduction never removes $x^*$: if the reduction to the bound $\beta_k$ of
coordinate $k$ is justified on a box $B\ni x^*$, then clamping $x^*_k$ to
$\beta_k$ gives a point of $B$ with value at most $\OPT$, which by uniqueness is
$x^*$; so $x^*_k=\beta_k$ and coordinate $k$ is active. Hence the current
face box always contains $x^*$, agrees with it on fixed coordinates, and has
$\norm{c-x^*}_\infty\le r$. For an active $i$ at $l_i$ that is not yet fixed,
$l_i$ is an endpoint of the current interval and
$\partial_iF(c)-C_2r\ge\gamma-2C_2r\ge\gamma/2>0$, so the midpoint test
fixes it; the upper case is symmetric. When no reduction applies, the free
set $J$ consists exactly of the continuous coordinates that are interior at
$x^*$. If $J=\emptyset$ the trial succeeds. Otherwise $x^*\pm tv\in X$ for
$v\in\R^J$ and small $t$, so \eqref{eq:growth} and Taylor's formula give
$\nabla_JF(x^*)=0$ and $H_{JJ}(x^*)\succeq2g\,\mathrm{Id}$. By the row-sum
bound, $H_{JJ}(c)\succeq(2g-C_3r)\mathrm{Id}$, and
\[
H_{JJ}(c)-(C_3r+\sigma_{\mu_*})\mathrm{Id}\succeq(2g-2C_3r-\sigma_{\mu_*})\mathrm{Id}
\succ0,
\]
because $2C_3r\le L/(4\cdot4^{\mu_*})\le g/32$ and
$\sigma_{\mu_*}=L/(4\cdot4^{\mu_*})\le g/32$. Exact symmetric elimination
detects this.

Through stage $j_*$, trial $\mu_*$ performs the table work of
\eqref{eq:dpwork} with the cap \eqref{eq:statecap} on numbers of bit length
$\poly(I+j_*+\mu_*K)$ at fixed $d$ (Section~\ref{sec:polynomial}), plus at most
$n^2$ sign tests and one elimination per stage. This is
$A_*=f_d(p,\kappa)\poly(I+B_\gamma)$ bit operations. Phase $q$ allots
$\sum_{\mu=2}^q2^{q-\mu}<2^{q-1}$ operations, and trial $\mu_*$ completes in
phase $q_*=\mu_*+\lceil\log_2A_*\rceil$. The total is at most
$2^{q_*}\le2\cdot2^{\mu_*}A_*=O(\sqrt\kappa\,A_*)$, plus clock overhead
polylogarithmic in this bound.
\end{proof}

The output is an exact implicit description, not a rational coordinate claim:
for $(x^2-\frac12)^2+y+xy^2$ on $[0,1]^2$ the face is $y=0$ and the optimizer
is $(1/\sqrt2,0)$. Weak reductions preserve the minimum but not necessarily
every optimizer, so without uniqueness the certificate proves global
optimality of $\hat x$, not uniqueness in $X$. Trials must be dovetailed: a
trial with $4^\mu<8\kappa$ may run forever without exceeding its cap, and
increasing $\mu$ once per unsuccessful stage would make the grading depend on
$\gamma$.

The margin term and strict complementarity are both needed for certificates
of this form, as the next two examples show.

\begin{proposition}[The margin cost is intrinsic to rectangular sign certificates]\label{prop:margin}
For $n\ge1$ let $X=[0,\frac12]^{n+2}$ with variables $x_0,\ldots,x_n,y$, and
\[
r_0=x_0-\tfrac14,\quad r_i=x_i-\tfrac14x_{i-1}^2,\quad
h=x_n-\tfrac18x_{n-1}^2,\quad
G_n=\sum_{i=0}^nr_i^2+y^2+3yh.
\]
Then $G_n$ has degree four, a path decomposition with bags of size at most
three, $L=\frac{35}{16}$, the unique optimizer $(a,0)$ with
$a_i=2^{-(4\cdot2^i-2)}$, $\OPT=0$, point growth with $g=\frac9{64}$
($\kappa=\frac{140}9$), and strict complementarity at $y$ with
$\gamma=\frac32a_n$. Every success of the procedure on $G_n$ produces a
rational number with reduced denominator at least $2^{4\cdot2^n-2}$.
Consequently a successful trial $\mu$ reaches a stage $j$ with
$j+\mu K\ge4\cdot2^n-2-\log_2\omega_0$, where $K$ is its cap and $\omega_0$ a
common denominator of the data.
\end{proposition}

\begin{proof}
The recursion $a_0=\frac14$, $a_i=\frac14a_{i-1}^2$ gives the stated $a_i$,
all in $(0,\frac12)$, and $r(a)=0$. With $e=x-a$, $r_0=e_0$ and
$r_i=e_i-\beta_ie_{i-1}$ with $\beta_i=(x_{i-1}+a_{i-1})/4\in[0,\frac14]$, so
$\norm r\ge\frac34\norm e$ and $F_n=\sum_ir_i^2\ge\frac9{16}\norm{x-a}^2$.
Since $h=\frac12(x_n+r_n)$,
\[
G_n=\tfrac14(F_n+y^2)+\tfrac34(F_n-r_n^2)+\tfrac34(r_n+y)^2+\tfrac32yx_n ,
\]
and every term is nonnegative on $X$. Thus
$G_n\ge\frac9{64}(\norm{x-a}^2+y^2)$, with equality to zero only at $(a,0)$.
The diagonal second derivatives are $2+\frac34x_i^2-x_{i+1}\le\frac{35}{16}$
for $i<n-1$, the same minus $\frac34y$ for $i=n-1$, and $2$ for $x_n$ and
$y$. The bags are $\{x_{i-1},x_i\}$ and $\{x_{n-1},x_n,y\}$. Finally
$\partial_yG_n(a,0)=3h(a)=3(a_n-\frac12a_n)$.

Consider a success. A zero-gap success outputs the unique optimizer, whose
coordinate $a_n=2^{-(4\cdot2^n-2)}$ has the stated denominator. Otherwise
the integer part is empty, and by the argument in the proof of
Theorem~\ref{thm:boundary}(b) no sound reduction fixes an $x$-coordinate,
because all are interior at the unique optimizer. If $y$ is not fixed, the
patch test needs $H_{JJ}\succ0$ at $(a,0)$, which lies in the face box; but
the principal minor on $\{x_n,y\}$ is
$\det\bigl(\begin{smallmatrix}2&3\\3&2\end{smallmatrix}\bigr)=-5$. So $y$ is
fixed, and not at its upper bound, since that would need
$\partial_yG_n\le0$ at $(a,0)$. So it is fixed by a sound certificate of
$\partial_yG_n\ge0$ on a retained box
$B=\prod_i[\alpha_i,\beta_i]\times[0,\beta_y]$ that contains $(a,0)$. At its
point with $y=0$, $x_n=\alpha_n$ and $x_{n-1}=\beta_{n-1}$,
$\partial_yG_n=3(\alpha_n-\beta_{n-1}^2/8)$, so
$\alpha_n\ge\beta_{n-1}^2/8\ge a_{n-1}^2/8=a_n/2>0$; also $\alpha_n\le a_n$.
A rational in $(0,2^{-(4\cdot2^n-2)}]$ has reduced denominator at least
$2^{4\cdot2^n-2}$. Box endpoints and grid nodes of stage $j$ in trial $\mu$
have denominators dividing $\omega_02^{j+\mu K}$, which proves the last claim.
\end{proof}

Here $I=O(n\log n)$ while every successful certificate of the above form
needs exponentially many bits in $n$, at fixed $d$, $p$ and $\kappa$. Thus
the $B_\gamma$ term of Theorem~\ref{thm:boundary} cannot be removed without
a different kind of certificate, for instance one that keeps the correlation
between $x_{n-1}$ and $x_n$, such as the displayed identity.

\begin{proposition}[Weak complementarity defeats sign and patch tests]\label{prop:weakcompl}
On $X=[\frac12,1]\times[0,1]^2$ let $u=x^2-\frac12$ and
\[
W=u^2+w_1^2+w_2^2+3w_1w_2+u(w_1-w_2)
=\tfrac12(u^2+w_1^2+w_2^2)+\tfrac12(u+w_1-w_2)^2+4w_1w_2 .
\]
Then $W$ has degree four and one bag of size three; its unique optimizer is
$x^*=(1/\sqrt2,0,0)$ with $\OPT=0$; point growth holds with $g=\frac12$ and
$L=12$ ($\kappa=24$); and $\nabla_wW(x^*)=0$, so both active coordinates are
weakly complementary. For every rational box $B\subseteq X$ containing $x^*$,
none of $\partial_xW$, $\partial_{w_1}W$, $\partial_{w_2}W$ has constant sign on
$B$, and $\nabla^2W(x^*)$ is not positive semidefinite. Hence no trial of the
procedure ever succeeds on $W$.
\end{proposition}

\begin{proof}
The identity is direct expansion. On $X$ it gives $W\ge\frac12(u^2+\norm w^2)$,
and $u^2=(x-x^*_1)^2(x+x^*_1)^2\ge(x-x^*_1)^2$ because $x+x^*_1>1$. So
$W\ge\frac12\norm{(x,w)-x^*}^2$, with zero only at $x^*$. Also
$\partial_x^2W=12x^2-2+2(w_1-w_2)\le12$ and $\partial_{w_i}^2W=2$. The
$x$-interval of $B$ contains the irrational $x^*_1$ in its interior. On the
segment $w=0$ inside $B$, $\partial_xW=4xu$, $\partial_{w_1}W=u$ and
$\partial_{w_2}W=-u$, and $u$ changes sign at $x^*_1$. So no coordinate can be
fixed by a sound test, and the patch test on all three coordinates would
need $\nabla^2W(x^*)\succ0$; but $v=(0,1,-1)$ gives
$v^{\mathsf T}\nabla^2W(x^*)v=2+2-6=-2$. A zero certified gap is impossible
because $x^*$ is irrational.
\end{proof}

The identity in Proposition~\ref{prop:weakcompl} is itself a short global
certificate that $\OPT=0$ and that $S=\{x^2=\frac12,\ w=0\}$. Certificates
that keep such correlations, rather than certify signs over independent
coordinate rectangles, are the natural route to exact output under point
growth alone; no general construction of them with parameterized size is
known.
